% =====================================================================
%  Agile & DevOps in Practice - Individual Final Assessment
%  Student Task Manager
%
%  Compiles with pdflatex (twice, for the table of contents):
%      pdflatex assessment-report.tex
%      pdflatex assessment-report.tex
%  Or open this file in Overleaf and press Recompile.
% =====================================================================
\documentclass[11pt,a4paper]{article}

\usepackage[T1]{fontenc}
\usepackage[utf8]{inputenc}
\usepackage{lmodern}
\usepackage[english]{babel}
\usepackage[a4paper,margin=2.4cm]{geometry}

\usepackage{booktabs}
\usepackage{tabularx}
\usepackage{longtable}
\usepackage{array}
\usepackage{enumitem}
\usepackage{graphicx}
\usepackage{xcolor}
\usepackage{listings}
\usepackage{caption}
\usepackage{titlesec}
\usepackage[hidelinks]{hyperref}

% ---------------------------------------------------------------------
%  Styling
% ---------------------------------------------------------------------
\definecolor{accent}{HTML}{1F4E79}
\definecolor{codebg}{HTML}{F5F6F8}
\definecolor{codeframe}{HTML}{D8DCE2}
\definecolor{good}{HTML}{1B6B3A}
\definecolor{bad}{HTML}{A22C22}

\hypersetup{
  colorlinks=true,
  linkcolor=accent,
  urlcolor=accent,
  citecolor=accent,
  pdftitle={Agile and DevOps in Practice - Individual Final Assessment},
  pdfauthor={Katy Great Adonai}
}

\titleformat{\section}{\color{accent}\Large\bfseries}{\thesection}{0.7em}{}
\titleformat{\subsection}{\color{accent}\large\bfseries}{\thesubsection}{0.7em}{}
\titleformat{\subsubsection}{\bfseries}{\thesubsubsection}{0.7em}{}

\captionsetup{font=small,labelfont={bf,color=accent}}

\newcolumntype{L}[1]{>{\raggedright\arraybackslash}p{#1}}

% Verbatim-style block for console output and code
\lstdefinestyle{plain}{
  basicstyle=\ttfamily\scriptsize,
  backgroundcolor=\color{codebg},
  frame=single,
  rulecolor=\color{codeframe},
  framesep=5pt,
  breaklines=true,
  columns=fullflexible,
  keepspaces=true,
  showstringspaces=false,
  literate={-}{-}1 {'}{'}1 {"}{"}1
}
\lstdefinestyle{yaml}{
  style=plain,
  language=,
  morecomment=[l]{\#},
  commentstyle=\color{accent!70}\itshape
}
\lstdefinestyle{java}{
  style=plain,
  language=Java,
  keywordstyle=\color{accent}\bfseries,
  commentstyle=\color{black!55}\itshape,
  stringstyle=\color{good}
}
\lstset{style=plain}

\newcommand{\yes}{\textcolor{good}{\textbf{Yes}}}
\newcommand{\no}{\textcolor{bad}{\textbf{No}}}
\newcommand{\code}[1]{\texttt{\small #1}}
\newcommand{\repo}{https://github.com/g1kdo/student-task-manager}

\setlist[itemize]{leftmargin=1.4em,itemsep=2pt,topsep=3pt}
\setlist[enumerate]{leftmargin=1.6em,itemsep=2pt,topsep=3pt}

% =====================================================================
\begin{document}

% ------------------------------ Title --------------------------------
\begin{titlepage}
\centering
\vspace*{2.5cm}

{\color{accent}\Huge\bfseries Agile \& DevOps in Practice\par}
\vspace{0.6cm}
{\LARGE Individual Final Assessment\par}

\vspace{1.6cm}
\rule{0.75\textwidth}{1pt}\par
\vspace{0.8cm}
{\LARGE\bfseries Project: Student Task Manager\par}
\vspace{0.4cm}
{\large A Java console application delivered over two simulated sprints\par}
\vspace{0.8cm}
\rule{0.75\textwidth}{1pt}\par

\vspace{2.2cm}
{\large\bfseries Katy Great Adona\"i\par}
\vspace{0.3cm}
{\large AmaliTech Training -- DEG Apprenticeship\par}

\vspace{1.8cm}
{\large\textbf{Repository:} \url{\repo}\par}

\vspace{0.5cm}
\begin{minipage}{0.82\textwidth}
\small\centering
Every claim in this report is checkable from the repository. Each user story
names the test that covers it, each review finding names the commit that closed
it, and the evidence files under \code{docs/evidence/} record the commands that
produced them. The build runs with a JDK alone:
\code{./mvnw clean verify}.
\end{minipage}

\vfill
{\large 4 September 2026\par}
\end{titlepage}

\tableofcontents
\newpage

% =====================================================================
\section{Overview}

\textbf{Student Task Manager} is a Java console application that lets a student
add tasks, list them with their status, mark them complete and delete them. It
is deliberately small: the assessment is about how the work was planned,
delivered, verified and reflected on, not about the size of the product.

This report consolidates every deliverable. The codebase, commit history,
pipeline definition and test suite are in the repository at
\url{\repo}.

\subsection{Submission history}

This is a resubmission. The first submission was reviewed and found to have
substantial gaps, chiefly in DevOps practice and reflection. Rather than present
a clean narrative that the commit history would contradict, this report:

\begin{itemize}
  \item records what each sprint actually delivered, including the pipeline
        deliverable that Sprint~1 did not really meet;
  \item states plainly that the sprint reviews and retrospectives were written
        after the review rather than during the sprints;
  \item maps each of the ten review findings to the commit that closed it
        (Section~\ref{sec:remediation}); and
  \item names the one finding that cannot be retrofitted --- the original
        single-day commit history --- instead of rewriting it.
\end{itemize}

\subsection{How to verify this report}

\begin{lstlisting}[style=plain]
git clone https://github.com/g1kdo/student-task-manager.git
cd student-task-manager
./mvnw clean verify        # compiles, runs 69 tests, coverage report + gate
./mvnw package -DskipTests
java -jar target/student-task-manager.jar
\end{lstlisting}

\noindent JDK~21 is the only requirement; the repository ships a Maven wrapper.
On Windows use \code{mvnw.cmd}.

% =====================================================================
\section{Sprint 0 --- Planning}

\subsection{Product vision}

\begin{quote}
A small Java console application that lets a student capture the tasks they have
to do, see what is still outstanding, and tick things off --- so that nothing is
tracked only in their head.
\end{quote}

\subsection{Product backlog}

Estimated in story points on a relative scale: 1 trivial, 2 small, 3 a
morning's work, 5 needs design as well as code.

\begin{table}[h]
\centering
\small
\begin{tabularx}{\textwidth}{@{}l X c c c@{}}
\toprule
\textbf{ID} & \textbf{User story} & \textbf{Priority} & \textbf{Points} & \textbf{Sprint} \\
\midrule
US1 & As a student, I want to add a task, so that I can stop keeping it in my head. & High & 3 & 1 \\
US2 & As a student, I want to see all my tasks with their status, so that I know what is outstanding. & High & 2 & 1 \\
US3 & As a student, I want to mark a task complete, so that I can track my progress. & High & 3 & 2 \\
US4 & As a student, I want to delete a task, so that my list stays relevant. & Medium & 2 & 2 \\
US5 & As a developer, I want application logs and a health check, so that I can see how the running system is behaving. & Medium & 5 & 2 \\
\midrule
US6 & As a student, I want mistyped input to be reported rather than crash or be silently misread, so that I can correct it and carry on. & High & 3 & Groomed in \\
US7 & As a developer, I want the tests to run automatically on every push, so that a regression is caught before review. & High & 3 & Groomed in \\
\bottomrule
\end{tabularx}
\caption{Product backlog, 21 points. US1--US5 (13 points) were the original
Sprint~0 backlog. US6 and US7 were added at grooming after the Sprint~1
retrospective.}
\end{table}

US5 was kept in the backlog from the start rather than added at the end, so
observability was planned work rather than an afterthought.

\subsection{Acceptance criteria}

Each criterion is written so a test can decide it. The covering test is named in
brackets --- this is what makes the criteria checkable rather than a matter of
opinion.

\subsubsection*{US1 --- Add a task}
\begin{itemize}
  \item Given a non-empty name, the task is stored and the application confirms
        it. \hfill \code{ConsoleAppTest.addingATaskConfirmsAndStores}
  \item Tasks appear in the order they were added. \hfill
        \code{TaskManagerTest.AddTask.addPreservesInsertionOrder}
  \item A blank or whitespace-only name is refused and nothing is stored.
        \hfill \code{TaskManagerTest.AddTask.addRejectsBlankNames}
  \item Surrounding whitespace is trimmed. \hfill
        \code{TaskManagerTest.AddTask.addTrimsTheName}
\end{itemize}

\subsubsection*{US2 --- View tasks}
\begin{itemize}
  \item Every task is listed, numbered from 1, showing \code{[Done]} or
        \code{[Pending]}, with a summary of how many of how many are complete.
        \hfill \code{ConsoleAppTest.listShowsNumbersAndStatuses}
  \item With no tasks, the application says so rather than printing an empty
        list or erroring. \hfill \code{ConsoleAppTest.viewingAnEmptyListSaysSo}
\end{itemize}

\subsubsection*{US3 --- Complete a task}
\begin{itemize}
  \item A valid task number becomes complete and is confirmed. \hfill
        \code{TaskManagerTest.CompleteTask.completeMarksTheTaskDone}
  \item Completing an already-complete task is reported and does not change the
        completed count. \hfill
        \code{TaskManagerTest.CompleteTask.completeIsIdempotent}
  \item A task number that does not exist --- including 0 and negatives --- is
        reported and leaves every task untouched. \hfill
        \code{completeRejectsOutOfRangeIndexes}, \code{taskNumberZeroIsReported}
  \item With no tasks, the application says so and does not ask for a number.
        \hfill \code{ConsoleAppTest.completingWithNoTasksSaysSo}
\end{itemize}

\subsubsection*{US4 --- Delete a task}
\begin{itemize}
  \item A valid task number is removed and confirmed; the remainder renumber
        without gaps. \hfill \code{deleteRemovesTheTask},
        \code{deletingRenumbersTheRemainder}
  \item Deleting a completed task lowers the completed count. \hfill
        \code{deleteUpdatesTheCompletedCount}
  \item A task number that does not exist is reported and nothing is removed.
        \hfill \code{deleteRejectsOutOfRangeIndexes}
\end{itemize}

\subsubsection*{US5 --- Logs and health check}
\begin{itemize}
  \item Every add, complete and delete produces a log record carrying a
        timestamp, a level and the source logger. \hfill
        \code{docs/evidence/demo-session.txt}
  \item A failed operation logs at \code{WARN}, not \code{INFO}. \hfill
        \code{docs/evidence/demo-session.txt}
  \item The health check reports live state --- uptime, task count, completed
        count, heap --- not a fixed string. \hfill
        \code{healthyApplicationReportsUp}, \code{healthCheckReportsLiveState}
  \item It reports \code{DEGRADED} under heap pressure and \code{DOWN} when the
        task store cannot be read, and never throws. \hfill
        \code{heapPressureReportsDegraded}, \code{unreadableTaskStoreReportsDown}
\end{itemize}

\subsubsection*{US6 --- Mistyped input}
\begin{itemize}
  \item A menu choice outside the six options is reported by name and the loop
        continues. \hfill \code{nonNumericChoiceIsReported},
        \code{outOfRangeChoiceIsReported}
  \item A task number that is not a whole number is reported and nothing
        changes. \hfill \code{nonNumericTaskNumberIsReported}
  \item Pressing Enter alone redraws the menu without reporting an error.
        \hfill \code{blankChoiceIsIgnored}
  \item Entering a task number never causes the \emph{next} menu read to be
        misread. \hfill two \code{REGRESSION}-named tests in
        \code{ConsoleAppTest}
  \item Reaching the end of input exits cleanly instead of looping forever.
        \hfill \code{endOfInputExitsCleanly}
\end{itemize}

\subsubsection*{US7 --- Automated build and test}
\begin{itemize}
  \item \code{./mvnw clean verify} compiles, runs every test and reports
        coverage on a machine with only a JDK. \hfill
        \code{docs/evidence/test-run.txt}
  \item The pipeline runs on every push and pull request to \code{main} and
        \code{develop}. \hfill \code{.github/workflows/ci.yml}
  \item A failing test fails the pipeline. \hfill
        \code{docs/evidence/pipeline-fails-on-failing-test.txt}
  \item The pipeline runs the packaged artefact and asserts on its output.
        \hfill the smoke-test step in \code{ci.yml}
  \item Test reports, coverage and the runnable jar are published as build
        artefacts. \hfill \code{ci.yml}
\end{itemize}

\subsection{Definition of Done}

The first version of this list --- \emph{code compiles, feature works, unit test
passes, code committed, no major bugs} --- could not fail anything, and did
not. Every item below can.

\begin{enumerate}
  \item \code{./mvnw clean verify} passes locally: compiles, all tests green,
        coverage gate met.
  \item Every acceptance criterion for the item is covered by at least one
        named automated test.
  \item Both the success \emph{and} the failure path of the item are tested,
        not the happy path alone.
  \item The operation logs at a level matching what happened --- \code{INFO}
        for normal work, \code{WARN} or \code{ERROR} for a problem.
  \item The work is committed in increments, each message saying what changed
        and why.
  \item The CI pipeline is green on the pushed commit.
  \item \code{README.md} and the relevant sprint document reflect any change in
        behaviour.
  \item No deliverables --- screenshots, reports --- are committed inside the
        source tree.
\end{enumerate}

\subsection{Sprint plans}

\textbf{Sprint 1:} US1 (3), US2 (2) --- 5 points, plus repository and pipeline
setup. Chosen because the list is useless until something can be put on it and
read back, and everything later depends on the storage model these two stories
establish.

\medskip
\noindent\textbf{Sprint 2:} US3 (3), US4 (2), US5 (5) --- 10 points. Completing
and deleting are what turn a list into a tracker; US5 makes the running
application observable.

% =====================================================================
\section{Sprint 1 --- Execution}

\subsection{Stories delivered}

5 of 5 committed points.

\begin{description}[leftmargin=0pt,style=nextline]
  \item[US1 --- Add a task] \code{TaskManager.addTask(String)} stores a
    \code{Task} and confirms it. \code{Task} is a small model holding a name and
    a completion flag. \\ Commits \code{a95258a}, \code{27fdf4c}, \code{8aeb39e}.
  \item[US2 --- View tasks] Tasks listed in insertion order, numbered from 1,
    each showing done or pending. \\ Commit \code{8aeb39e}.
\end{description}

\subsection{Repository and pipeline setup}

\begin{itemize}
  \item Git repository with \code{main} and \code{develop} branches and a pull
        request template.
  \item A GitHub Actions workflow checking out the code and installing
        Temurin~21 (\code{629bd17}).
  \item Two JUnit~5 tests for \code{addTask} and \code{completeTask}
        (\code{88e15bb}).
\end{itemize}

\subsection{A real problem diagnosed}

\code{Task} was first placed in a package named \code{java}. The JVM refuses to
load user classes from a package under \code{java.}, so the class would not run.
The package was renamed to \code{app}; commit \code{27fdf4c} records the
diagnosis, not merely the change.

\subsection{Sprint 1 review: goal met, with a significant caveat}

Both stories work and satisfy their acceptance criteria. The \emph{pipeline}
deliverable, however, was not met, and this was not recognised at the time:

\begin{itemize}
  \item The workflow's only test step was committed as a comment
        (\code{729b878}, ``Comment out Maven test command in workflow''). The
        pipeline could not fail, so it verified nothing.
  \item The step was commented out because \code{mvn test} did not work --- and
        \code{mvn test} did not work because the repository had no
        \code{pom.xml} or \code{build.gradle} at all. The two JUnit tests were
        only runnable inside one IntelliJ project.
  \item Commenting out the check treated the symptom. The cause was a missing
        build file.
\end{itemize}

\subsection{Sprint 1 retrospective}

\noindent\textit{No retrospective was written during Sprint~1. That omission is
itself the sprint's clearest process failure, and it is recorded here rather
than disguised; this reflection was written on 4 September 2026.}

\subsubsection*{What went well}
\begin{itemize}
  \item Both committed stories were delivered and worked.
  \item Observability was planned, not bolted on: US5 sat in the backlog from
        Sprint~0.
  \item Bounds checks were written deliberately, so an out-of-range task number
        printed a message instead of throwing.
  \item One real diagnosis was made and recorded in its commit message.
\end{itemize}

\subsubsection*{What was difficult}
\begin{itemize}
  \item \textbf{The build tool was the thing I did not know I needed.} I had
        JUnit tests and a CI workflow and assumed that was a pipeline. Without
        a \code{pom.xml}, \code{mvn test} had nothing to run.
  \item \textbf{I treated a red pipeline as a problem with the pipeline.} When
        \code{mvn test} failed I commented the step out, which made the symptom
        go away and left the cause in place. A pipeline that cannot fail is not
        a pipeline, and I had built exactly that.
  \item \textbf{I tested what I expected to happen.} Two tests, both happy
        paths. I never tested a task number that does not exist, even though I
        had written the guard for it.
  \item \textbf{I could not separate the front end from the rules.}
        \code{TaskManager} printed to \code{System.out}, so its behaviour could
        only be checked by looking at the console. That is why the menu loop
        went untested, and why a user-visible bug in it shipped.
  \item \textbf{I committed in one sitting.} Nine of ten commits landed on a
        single day.
\end{itemize}

\subsubsection*{Improvements for Sprint 2}

\begin{table}[h]
\centering
\small
\begin{tabularx}{\textwidth}{@{}c X L{5.1cm}@{}}
\toprule
\textbf{\#} & \textbf{Improvement} & \textbf{How it is verified} \\
\midrule
1 & Add a real build file declaring the JUnit dependency, and make CI run \code{mvn test} instead of commenting it out. & \code{docs/evidence/test-run.txt} \\
2 & Never disable a check to make a build green; fix the cause and prove the pipeline can fail. & \code{docs/evidence/pipeline-\allowbreak fails-on-failing-test.txt} \\
3 & Test the failure path of every story, starting with the out-of-range task numbers whose guards were already written. & Out-of-range, blank-name and repeat-completion tests \\
4 & Get printing out of \code{TaskManager} so the rules and the menu can both be tested. & \code{ConsoleApp} takes its streams as arguments; 19 menu tests \\
5 & Write the review and retrospective as the sprint's work, not after it. & \code{sprint2-review.md}, \code{sprint2-retrospective.md} \\
\bottomrule
\end{tabularx}
\caption{Sprint 1 improvement actions, each paired with the artefact that
proves it was carried out.}
\end{table}

Two of these were turned into tracked backlog items rather than handled
informally: \textbf{US7} (build and test on every push) from improvements 1--2,
and \textbf{US6} (mistyped input) from improvements 3--4.

% =====================================================================
\section{Sprint 2 --- Execution and Improvement}

\subsection{Stories delivered}

10 of 10 committed points.

\subsubsection*{US3 --- Complete a task}

\code{TaskManager.completeTask(int)} marks the task at a zero-based index
complete and \emph{reports whether it changed anything}, so completing an
already-complete task is distinguishable from a real change and cannot inflate
the completed count. An index outside the list is refused and logged at
\code{WARN}. Nine tests in \code{TaskManagerTest.CompleteTask}, plus menu-level
coverage.

\subsubsection*{US4 --- Delete a task}

\code{TaskManager.deleteTask(int)} removes the task at an index, renumbers the
remainder by list order, and keeps the completed count correct when the deleted
task was complete. Nine tests in \code{TaskManagerTest.DeleteTask}.

\subsubsection*{US5 --- Logging}

SLF4J with a Logback backend, configured in
\code{src/main/resources/logback.xml}. Every record carries a timestamp, a level
and the source logger, and goes to both the console and a daily rolling file
under \code{logs/} kept for seven days. Levels are used as levels:

\begin{itemize}
  \item \code{INFO} --- normal work: task added, completed, deleted;
        application starting and shutting down.
  \item \code{WARN} --- a recoverable problem: unknown task number, blank name,
        unrecognised menu choice.
  \item \code{ERROR} --- a failure: console input failing, health check
        reporting \code{DOWN}.
\end{itemize}

The first attempt's ``logging'' was two \code{System.out.println("[LOG] ...")}
lines in the add and view branches --- no framework, no levels, no timestamps,
and nothing at all for complete or delete. All four operations now log, and they
log from \code{TaskManager} where the state actually changes rather than from
the menu.

\subsubsection*{US5 --- Health check}

\code{HealthCheck.check()} inspects live state and returns a
\code{HealthStatus}:

\begin{table}[h]
\centering
\small
\begin{tabularx}{\textwidth}{@{}l X l@{}}
\toprule
\textbf{Check} & \textbf{What it does} & \textbf{Verdict it can produce} \\
\midrule
\code{taskStore} & Reads the task count and completed count & \code{DOWN} if reading throws \\
\code{heap} & Compares heap in use against the maximum & \code{DEGRADED} at or above 90\% \\
\bottomrule
\end{tabularx}
\caption{The two checks and the bad news each can report.}
\end{table}

It also reports uptime, task count, completed count and heap figures, and never
throws --- a health check that fails by throwing cannot report that the
application is unwell. The verdict is logged at a level matching its severity.
The first attempt's health check printed the fixed string
\code{"Application is running."}, which was true by construction and inspected
nothing.

The clock and the heap probe are constructor arguments, so \code{DEGRADED} and
\code{DOWN} are asserted directly rather than hoped for:

\begin{lstlisting}[style=java]
@Test
@DisplayName("an unreadable task store reports DOWN without throwing")
void unreadableTaskStoreReportsDown() {
    TaskManager broken = new TaskManager() {
        @Override public int getTaskCount() {
            throw new IllegalStateException("task store unavailable");
        }
    };

    HealthStatus health =
            new HealthCheck(broken, new TickingClock(), heapAt(0, ONE_GIB)).check();

    assertEquals(HealthStatus.Status.DOWN, health.status());
    assertEquals(-1, health.taskCount(), "an unreadable store reports -1, not a made-up 0");
}
\end{lstlisting}

\subsection{Sprint 2 review: goal met}

All three stories work, each acceptance criterion is covered by a named test,
and the application is observable through both a log trail and a health check
that can actually report bad news. The demonstration is in
Section~\ref{sec:demo}.

\subsection{Improvements from Sprint 1 applied}

\begin{table}[h]
\centering
\small
\begin{tabularx}{\textwidth}{@{}c X L{4.6cm}@{}}
\toprule
\textbf{\#} & \textbf{Improvement} & \textbf{Status} \\
\midrule
1 & Add a real build file and make CI run the tests & Done --- \code{ac292c5}, \code{1aa88a8} \\
2 & Never disable a check to go green; prove the pipeline can fail & Done --- evidence recorded \\
3 & Test the failure path of every story & Done --- 69 tests, 92.3\% lines \\
4 & Get printing out of \code{TaskManager} & Done --- \code{f6a9182} \\
5 & Write the review and retrospective as the sprint's work & Done --- this section \\
\bottomrule
\end{tabularx}
\caption{Sprint 1 improvements, and where each was carried out.}
\end{table}

\subsection{Sprint 2 retrospective}

\subsubsection*{What went well}
\begin{itemize}
  \item All three committed stories were delivered, and US5 went from a
        placeholder to real behaviour: levelled logging across all four
        operations, and a health check with three possible verdicts rather than
        one.
  \item \textbf{The menu input validation held.} Checking
        \code{choice.matches("\textbackslash\textbackslash d+")} before
        \code{Integer.parseInt} was deliberate and prevented the
        \code{InputMismatchException} crash a \code{Scanner} menu loop normally
        suffers. That instinct was right; it was simply applied in one place and
        not the other.
  \item \textbf{Separating rendering from the rules paid for itself
        immediately.} It was done to make the rules testable; it made the
        \emph{menu} testable too, and that is what caught the bug below.
  \item \textbf{Making the health check injectable was worth the extra
        interface.} Testing an unhappy path you cannot trigger is not testing
        it.
\end{itemize}

\subsubsection*{What was difficult}
\begin{itemize}
  \item \textbf{A user-visible bug shipped, and it shipped precisely where
        nothing was tested.} Menu options 3 and 4 read the task number with
        \code{scanner.nextInt()} inside a loop whose next read was
        \code{scanner.nextLine()}. \code{nextInt()} consumes the digits but
        leaves the newline, so the following \code{nextLine()} returned an empty
        string and the menu printed ``Input a valid digit'' after every complete
        and every delete. The lesson is not about \code{Scanner}: the bug was in
        the one layer that had no tests, and it survived because I checked the
        code by reading it rather than by running it. Two tests named
        \code{REGRESSION} now pin the fix.
  \item \textbf{I mistook configuring a thing for having a working thing.} The
        same mistake as Sprint~1's commented-out test step, in a different
        place: two \code{[LOG]} printlns looked like logging and a fixed status
        string looked like a health check. Both existed; neither did anything.
        The question I should have asked of each is: \emph{what would this tell
        me if something were wrong?}
  \item \textbf{I mixed deliverables into the source tree.} Screenshots and the
        report PDF were committed under \code{src/app/screenshots/}.
  \item \textbf{The two sprints were one day.} Nine of ten commits landed on
        6~July 2026. Whatever the commit messages say, the history does not show
        two sprints of iterative delivery, and I cannot retrofit that.
\end{itemize}

\subsubsection*{Lessons across both sprints}

\begin{enumerate}
  \item \textbf{A check that cannot fail is not a check.} This covers the
        commented-out test step, the \code{[LOG]} printlns and the constant
        health string alike. In all three cases something existed where a check
        was expected, which is worse than nothing, because it stops you looking.
        Each is now something that can report bad news --- and the pipeline
        failing on a deliberately inverted assertion
        (Section~\ref{sec:canfail}) is the proof.
  \item \textbf{Testability is a design property, not a testing activity.} I
        could not test the menu because \code{ConsoleApp} did not exist and
        \code{TaskManager} printed to a global stream. The tests became possible
        once the design let them be written --- the streams became arguments,
        the clock became an argument. Coverage went from 2 tests to 69 without
        the rules themselves getting more complicated.
  \item \textbf{Bugs cluster where the tests are not.} The one user-visible
        defect was in the only untested layer. That is not coincidence, and it
        is the most useful thing I learned.
  \item \textbf{When a build goes red, the build is usually telling the truth.}
        My instinct was to silence it. The red step was correctly reporting that
        there was no build file.
\end{enumerate}

\subsubsection*{What I would do differently from day one}
\begin{itemize}
  \item Create the \code{pom.xml} in the first commit, before any code, so the
        tests are runnable from outside the IDE from the beginning.
  \item Write the failing test for a story's error path before the code that
        handles it, rather than writing a guard and never checking it.
  \item Take the streams and the clock as constructor arguments from the start.
  \item Commit the review and retrospective as part of the sprint, on the day,
        so the reflection is a deliverable rather than a reconstruction.
  \item Keep \code{docs/} and \code{src/} separate from the first commit.
\end{itemize}

% =====================================================================
\section{Remediation --- responding to the assessment review}
\label{sec:remediation}

Each finding from the review, the change that closed it, and where to check.

\begin{longtable}{@{}c L{4.5cm} L{5.3cm} l L{2.7cm}@{}}
\toprule
\textbf{\#} & \textbf{Finding} & \textbf{Change} & \textbf{Commit} & \textbf{Evidence} \\
\midrule
\endfirsthead
\toprule
\textbf{\#} & \textbf{Finding} & \textbf{Change} & \textbf{Commit} & \textbf{Evidence} \\
\midrule
\endhead
\bottomrule
\endfoot
1 & No build file anywhere, so \code{mvn test} could never work and the tests were unrunnable outside IntelliJ & Maven build targeting Java 21 declaring \code{junit-jupiter}, surefire failing the build on a failing test, plus a wrapper so it runs on a clean machine & \code{ac292c5} & \code{test-run.txt} \\
\addlinespace
2 & The only test step was commented out; a pipeline that cannot fail is not a pipeline & Workflow compiles, runs the suite, packages the jar, produces coverage and smoke-tests the artefact; triggers cover \code{develop} too & \code{1aa88a8} & \code{pipeline-fails-\allowbreak on-failing-test.txt} \\
\addlinespace
3 & ``Logging'' was two hardcoded printlns, no levels or timestamps, nothing for complete or delete & SLF4J + Logback; all four operations log with timestamps and levels; \code{WARN}/\code{ERROR} used meaningfully; console plus rolling file & \code{f6a9182} & \code{demo-session.txt} \\
\addlinespace
4 & The health check printed a fixed string and inspected nothing & Reads the task store, measures heap against a 90\% threshold, reports uptime and counts; \code{UP}/\code{DEGRADED}/\code{DOWN}; never throws & \code{57478b0} & \code{HealthCheckTest} \\
\addlinespace
5 & \code{nextInt()} then \code{nextLine()} made ``Input a valid digit'' print after every complete and delete & All input read by line through one prompt; rendering moved into a new \code{ConsoleApp} taking its streams as arguments & \code{f6a9182} & two \code{REGRESSION} tests \\
\addlinespace
6 & Two tests, happy paths of add and complete only & 69 tests across four classes; out-of-range indexes parametrically including \code{MIN\_VALUE}/\code{MAX\_VALUE}; blank names, repeat completion, end of input & \code{dcd4276}, \code{f9aa044}, \code{1de2787} & \code{test-run.txt}, \code{coverage.txt} \\
\addlinespace
7 & No sprint review and no retrospective anywhere & Four documents written, with the timing stated honestly rather than backdated & \code{b858d1c}, \code{f40782d} & \code{docs/sprint*.md} \\
\addlinespace
8 & Screenshots and a PDF committed inside \code{src/app/} & Restructured to the Maven standard layout; deliverables moved to \code{docs/}; the IntelliJ \code{.iml} dropped & \code{6962ae7} & repository tree \\
\addlinespace
9 & The one later commit added an unrelated \code{StreamArrayExample} & Moved to an \code{examples} package so \code{app} holds only assessment code, and excluded from the coverage gate & \code{6962ae7} & \code{src/main/java/\allowbreak examples/} \\
\addlinespace
10 & Nine of ten commits on one day, so both sprints were a single-day exercise & Not retrofittable, and not disguised. The remediation is 12 scoped commits; the original history is left intact. & --- & \code{git log} \\
\end{longtable}

\subsection{Before and after}

\begin{table}[h]
\centering
\small
\begin{tabularx}{\textwidth}{@{}X l l@{}}
\toprule
& \textbf{First submission} & \textbf{Now} \\
\midrule
Tests & 2 & 69 \\
Runnable outside the IDE & \no & \yes\ (\code{./mvnw test}) \\
Line coverage & Not measurable & 92.3\% \\
Branch coverage & Not measurable & 89.2\% \\
Pipeline can fail & \no & \yes, demonstrated \\
Operations that log & 2 of 4, no levels & 4 of 4, levels + timestamps \\
Health check verdicts & 1 (a constant) & 3 (\code{UP}, \code{DEGRADED}, \code{DOWN}) \\
Sprint reviews / retrospectives & 0 / 0 & 2 / 2 \\
\bottomrule
\end{tabularx}
\caption{What changed, measured rather than asserted.}
\end{table}

% =====================================================================
\section{DevOps --- the CI/CD pipeline}

\subsection{The build}

\code{pom.xml} targets Java 21 and declares \code{junit-jupiter} for tests and
SLF4J with Logback for logging. Four plugins do the work: \textbf{compiler},
\textbf{surefire} (a failing test fails the build), \textbf{JaCoCo} (coverage
report, and a gate that fails \code{verify} below 85\% line or 80\% branch
coverage) and \textbf{jar} (a runnable jar with \code{app.Main} as its entry
point, with runtime dependencies copied alongside).

A script-only Maven wrapper is committed, so the build runs on a machine with
nothing but a JDK --- which is exactly the situation the CI runner is in, and
was the situation the first submission could not handle.

\subsection{The pipeline}

\code{.github/workflows/ci.yml} runs on every push and pull request to
\code{main} and \code{develop}, and can be re-run by hand.

\begin{table}[h]
\centering
\small
\begin{tabularx}{\textwidth}{@{}l X@{}}
\toprule
\textbf{Step} & \textbf{What it does, and what it would catch} \\
\midrule
Checkout & Fetches the commit under test \\
Set up JDK 21 & Temurin 21, with the Maven cache restored \\
Compile & \code{./mvnw clean compile} --- catches anything that does not build \\
Test & \code{./mvnw test} --- 69 tests; a failure fails the pipeline \\
Package + coverage & \code{./mvnw verify} --- builds the jar, writes the JaCoCo report, fails below the coverage floor \\
Smoke test & Drives the \emph{packaged jar} over stdin and greps its output for the confirmations and \code{status=UP}, and for the absence of the spurious rejection message --- catches a jar that compiles but does not run \\
Publish artefacts & Uploads surefire reports, the coverage report and the runnable jar \\
\bottomrule
\end{tabularx}
\caption{Pipeline steps. The smoke test is what makes this a check on the
artefact rather than only on the source.}
\end{table}

The smoke-test step, in full:

\begin{lstlisting}[style=yaml]
- name: Smoke-test the packaged application
  run: |
    printf '1\nWrite sprint report\n2\n3\n1\n2\n5\n6\n' \
      | java -jar target/student-task-manager.jar > smoke.out 2>&1
    fail() { echo "::error::$1"; exit 1; }
    grep -q 'Task added'          smoke.out || fail 'add task did not confirm'
    grep -q 'Write sprint report' smoke.out || fail 'list did not render the task'
    grep -q 'Task completed'      smoke.out || fail 'complete did not confirm'
    grep -q 'status=UP'           smoke.out || fail 'health check did not report UP'
    if grep -qi 'not a valid choice' smoke.out; then
      fail 'the menu rejected valid input (stale newline regression)'
    fi
\end{lstlisting}

\subsection{Evidence that the pipeline can fail}
\label{sec:canfail}

The central criticism of the first submission was that its pipeline could not
fail. Rather than assert that this one can, it was demonstrated: one assertion
in \code{TaskTest} was deliberately inverted, \code{./mvnw clean verify} was
run, and the result recorded before the assertion was restored.

\begin{lstlisting}[style=plain]
[ERROR] Tests run: 11, Failures: 1, Errors: 0 <<< FAILURE! -- in app.TaskTest
[ERROR] app.TaskTest.completeReportsTheChange <<< FAILURE!
org.opentest4j.AssertionFailedError: DELIBERATELY WRONG: proving the build
  fails on a failing test ==> expected: <false> but was: <true>
        at app.TaskTest.completeReportsTheChange(TaskTest.java:52)
[ERROR] Tests run: 69, Failures: 1, Errors: 0, Skipped: 0
[INFO] BUILD FAILURE
\end{lstlisting}

\noindent Recorded in \code{docs/evidence/pipeline-fails-on-failing-test.txt},
which also notes that restoring the assertion returned the build to
\code{BUILD SUCCESS}.

\subsection{Pipeline runs}

The GitHub Actions run for the current pipeline is at
\url{\repo/actions}. For comparison, the first attempt's run --- checkout and
setup-java only, with no test step --- is archived at
\code{docs/archive/first-attempt/cicd1-no-test-step.png}.

% Replace with a screenshot of the current green run once pushed:
% \begin{figure}[h]
%   \centering
%   \includegraphics[width=0.95\textwidth]{../evidence/ci-run.png}
%   \caption{The CI pipeline running green: compile, test, package, smoke test.}
% \end{figure}

% =====================================================================
\section{Testing evidence}

\subsection{The suite}

\begin{table}[h]
\centering
\small
\begin{tabularx}{\textwidth}{@{}l c X@{}}
\toprule
\textbf{Test class} & \textbf{Tests} & \textbf{What it covers} \\
\midrule
\code{TaskManagerTest} & 30 & The rules, grouped by story: insertion order and name trimming (US1); empty-list and completion-state reporting and the unmodifiable view (US2); idempotent completion and boundary indexes (US3); index shifting and completed-count bookkeeping (US4) \\
\code{TaskTest} & 11 & The model and the invariant it guards: blank and null names refused, trimming, idempotent completion, \code{toString} \\
\code{HealthCheckTest} & 9 & \code{UP} with correct counts, uptime from construction, \code{DEGRADED} at the threshold and \code{UP} just below it, \code{DOWN} on an unreadable store, and that heap pressure does not soften a \code{DOWN} verdict \\
\code{ConsoleAppTest} & 19 & Integration: the real menu loop driven by scripted keystrokes. Two \code{REGRESSION} tests for the reported input bug, the full menu contract, empty names, task number 0, non-numeric task numbers, renumbering after a delete, end of input, and one full session across all five stories \\
\midrule
\textbf{Total} & \textbf{69} & \\
\bottomrule
\end{tabularx}
\caption{Test suite, up from two happy-path tests.}
\end{table}

Out-of-range indexes are checked parametrically, including
\code{Integer.MIN\_VALUE} and \code{Integer.MAX\_VALUE}. Each
\code{ConsoleAppTest} session runs under a timeout, so a menu loop that
mishandles end of input fails the build instead of hanging it.

\subsection{Result}

\begin{lstlisting}[style=plain]
[INFO] Running app.ConsoleAppTest
[INFO] Tests run: 19, Failures: 0, Errors: 0, Skipped: 0
[INFO] Running app.HealthCheckTest
[INFO] Tests run: 9, Failures: 0, Errors: 0, Skipped: 0
[INFO] Running app.TaskManagerTest
[INFO] Tests run: 30, Failures: 0, Errors: 0, Skipped: 0
[INFO] Running app.TaskTest
[INFO] Tests run: 11, Failures: 0, Errors: 0, Skipped: 0
[INFO] Tests run: 69, Failures: 0, Errors: 0, Skipped: 0
[INFO] All coverage checks have been met.
[INFO] BUILD SUCCESS
\end{lstlisting}

\subsection{Coverage}

\begin{table}[h]
\centering
\small
\begin{tabular}{@{}l r r r@{}}
\toprule
\textbf{Class} & \textbf{Instructions} & \textbf{Branches} & \textbf{Lines} \\
\midrule
\code{TaskManager}  & 100.0\% & 100.0\% & 31/31 \\
\code{Task}         & 100.0\% & 100.0\% & 15/15 \\
\code{HealthCheck}  & 100.0\% &  90.0\% & 39/39 \\
\code{HealthStatus} & 100.0\% & 100.0\% & 12/12 \\
\code{ConsoleApp}   &  87.9\% &  83.3\% & 102/119 \\
\midrule
\textbf{Total}      & \textbf{94.5\%} & \textbf{89.2\%} & \textbf{205/222 (92.3\%)} \\
\bottomrule
\end{tabular}
\caption{JaCoCo coverage. The \code{examples} package and the \code{Main} entry
point are excluded; \code{Main} is covered by the CI smoke test instead. The
build fails below 85\% line or 80\% branch coverage.}
\end{table}

% =====================================================================
\section{Prototype demonstration}
\label{sec:demo}

A complete session, reproduced with
\code{java -jar target/student-task-manager.jar} and recorded in
\code{docs/evidence/demo-session.txt}. Menu redraws are elided; log lines are
shown as the application emits them.

\begin{lstlisting}[style=plain]
11:04:33.876 INFO  ConsoleApp  - Student Task Manager starting up

Enter your choice: 1
Task name: Read chapter 1
11:04:33.876 INFO  TaskManager - Task added: 'Read chapter 1' (total now 1)
Task added.

Enter your choice: 1
Task name: Write summary
11:04:33.878 INFO  TaskManager - Task added: 'Write summary' (total now 2)
Task added.

Enter your choice: 2

Your tasks:
  1. Read chapter 1                 [Pending]
  2. Write summary                  [Pending]
  0 of 2 complete.

Enter your choice: 3
Task number: 2
11:04:33.890 INFO  TaskManager - Task completed: 'Write summary' (1 of 2 now done)
Task completed.

Enter your choice: 4
Task number: 9
11:04:33.894 WARN  TaskManager - Cannot delete task 9: no such task (there are 2)
Task 9 was not found.

Enter your choice: 4
Task number: 1
11:04:33.895 INFO  TaskManager - Task deleted: 'Read chapter 1' (1 remaining)
Task deleted.

Enter your choice: 2

Your tasks:
  1. Write summary                  [Done]
  1 of 1 complete.

Enter your choice: 5
11:04:33.893 INFO  HealthCheck - Health check: status=UP uptime=0h00m00s
                   tasks=1 completed=1 heap=12.9MB/4006.0MB (0.3%)

Health: status=UP uptime=0h00m00s tasks=1 completed=1 heap=12.9MB/4006.0MB (0.3%)
  taskStore  readable, holding 1 task(s)
  heap       0.3% used

Enter your choice: x
'x' is not a valid choice. Enter a number from 1 to 6.
11:04:33.895 WARN  ConsoleApp  - Unrecognised menu choice: 'x'

Enter your choice: 6
Goodbye.
11:04:33.896 INFO  ConsoleApp  - Student Task Manager shut down
\end{lstlisting}

Three things in this transcript are worth noting against the review findings.
Option~3 is followed immediately by option~4 with no spurious rejection in
between, which is the reported bug's absence. The out-of-range delete logs at
\code{WARN} rather than \code{INFO}. And the health check reports
\code{tasks=1 completed=1}, following the real state after the delete rather
than printing a constant.

\subsection{Architecture}

Five small classes with one job each, which is what makes the tests possible.

\begin{table}[h]
\centering
\small
\begin{tabularx}{\textwidth}{@{}l X@{}}
\toprule
\textbf{Class} & \textbf{Responsibility} \\
\midrule
\code{Task} & The model. Guards its own invariant: a blank name is refused at construction \\
\code{TaskManager} & The rules. Owns the list, reports success through return values, logs what happened. Never writes to \code{System.out} \\
\code{HealthCheck} / \code{HealthStatus} & Observability. Takes a clock and a heap probe as arguments so the unhappy paths are testable \\
\code{ConsoleApp} & The front end. Owns all rendering; takes its streams as arguments so tests drive the real menu \\
\code{Main} & The wiring that supplies the real console \\
\bottomrule
\end{tabularx}
\end{table}

% =====================================================================
\section{Delivery discipline}

The current commit history, oldest first:

\begin{table}[h]
\centering
\small
\begin{tabularx}{\textwidth}{@{}l l X@{}}
\toprule
\textbf{Commit} & \textbf{Date} & \textbf{Subject} \\
\midrule
\code{a95258a} & 2026-07-06 & First commit: initial project setup \\
\code{27fdf4c} & 2026-07-06 & Moved Task into app/ since java/ is a reserved package \\
\code{8aeb39e} & 2026-07-06 & Implemented the TaskManager \\
\code{88e15bb} & 2026-07-06 & Tested the codes using unit testing \\
\code{629bd17} & 2026-07-06 & Creating workflow... \\
\code{729b878} & 2026-07-06 & Comment out Maven test command in workflow \\
\code{18ab186} & 2026-07-06 & Implemented the main console UI and documentation \\
\code{99cb8f9} & 2026-07-06 & Screenshot of the last build attempt \\
\code{7e36ea8} & 2026-07-23 & Update documentation and add stream example \\
\midrule
\code{6962ae7} & 2026-09-04 & chore: restructure repository into Maven standard layout \\
\code{ac292c5} & 2026-09-04 & build: define a Maven build so the JUnit tests are runnable \\
\code{1aa88a8} & 2026-09-04 & ci: run the tests in the pipeline instead of commenting them out \\
\code{f6a9182} & 2026-09-04 & fix: read all console input by line, split rendering from the domain \\
\code{57478b0} & 2026-09-04 & feat: make the health check inspect real state \\
\code{dcd4276} & 2026-09-04 & test: cover the delete, invalid-index and blank-name paths \\
\code{f9aa044} & 2026-09-04 & test: cover the health check, including DEGRADED and DOWN \\
\code{1de2787} & 2026-09-04 & test: drive the whole console menu, pinning the reported bug \\
\code{fc7e905} & 2026-09-04 & docs: capture reproducible test, coverage and demo evidence \\
\code{b858d1c} & 2026-09-04 & docs: add the Sprint 1 review and retrospective \\
\code{f40782d} & 2026-09-04 & docs: add the Sprint 2 review, retrospective and remediation record \\
\code{53ab06b} & 2026-09-04 & docs: rewrite the README around the real build and deliverables \\
\bottomrule
\end{tabularx}
\caption{Commit history. The rule above the last block separates the original
submission from the remediation work.}
\end{table}

The original nine commits do describe real ordering --- setup, then
\code{TaskManager}, then tests, then the workflow, then the console UI --- and
one of them records a genuine diagnosis. But they landed in a single day, so the
history does not show two sprints of iterative delivery. That is stated here
rather than concealed, and the original commits were left as they are rather
than rewritten.

The remediation work is twelve commits, each one scoped change with a message
saying what changed and why. Every one of them builds and every one leaves the
test suite green.

% =====================================================================
\section{Summary of deliverables}

\begin{table}[h]
\centering
\small
\begin{tabularx}{\textwidth}{@{}L{4.4cm} X@{}}
\toprule
\textbf{Deliverable} & \textbf{Where} \\
\midrule
Backlog, estimates, acceptance criteria, Definition of Done, sprint plans
  & Section 2 of this report; \code{docs/sprint0-planning.md} \\
Codebase with incremental commit history
  & \url{\repo}; Section 8 \\
CI/CD evidence
  & \code{.github/workflows/ci.yml}; Section 6;
    \code{docs/evidence/pipeline-fails-on-failing-test.txt} \\
Testing evidence
  & 69 tests in \code{src/test/java/app/}; Section 7;
    \code{docs/evidence/test-run.txt}, \code{coverage.txt} \\
Sprint review documents
  & Sections 3.4 and 4.2; \code{docs/sprint1-review.md},
    \code{docs/sprint2-review.md} \\
Retrospectives
  & Sections 3.5 and 4.4; \code{docs/sprint1-retrospective.md},
    \code{docs/sprint2-retrospective.md} \\
Response to the assessment review
  & Section 5; \code{docs/remediation-sprint.md} \\
Prototype demonstration
  & Section 9; \code{docs/evidence/demo-session.txt} \\
First attempt, kept for comparison
  & \code{docs/archive/} \\
\bottomrule
\end{tabularx}
\end{table}

\vfill
\begin{center}
\small
\textit{Katy Great Adona\"i --- AmaliTech Training, DEG Apprenticeship ---
4 September 2026}
\end{center}

\end{document}
